\chapter{Introduction}\label{ch:introduction}

\section{Motivation}\label{sec:motivation}

Modern software is assembled from external libraries, called \emph{dependencies}. Each library exists in many versions, and each version expects particular versions of other libraries and of the programming language itself. Together, these requirements form the \emph{dependency configuration} of a project. When this configuration is missing, incomplete or out of date, the project fails to run even though its own source code is correct.

The problem is especially damaging for shared scientific code. Researchers increasingly publish computational notebooks, typically Jupyter notebooks, together with the paper that reports their results, so that others can re-run the analysis. Libraries evolve, however: old versions become uninstallable, functions are removed, and the interpreter itself moves on. A notebook that ran on the day of publication can therefore stop executing years later, and with it the possibility of confirming the reported results. Empirical studies show that this is not a rare corner case. Pimentel et al.\ examined notebooks mined from GitHub and found that only a minority of them can be re-executed and reproduce their results \cite{pimentel2019large}. Samuel and Mietchen re-executed notebooks that accompany biomedical publications and report on how many of them fail and why \cite{samuel2024computational}.\todo{Quote the headline reproducibility numbers of both studies; I have not re-read them for this draft.} % TODO: author -- quote the headline numbers of [pimentel2019large] and [samuel2024computational] after re-reading them; they are deliberately not given from memory.
Missing or broken dependencies are among the leading causes of these failures \cite{samuel2024computational}.% TODO: author to confirm -- statement taken from the vision document; check against the paper.
Jin et al.\ report the same erosion for machine-learning notebooks: of 12{,}106 notebooks from 75 Kaggle competitions, only 26\,\% remain reproducible today. They also observe that rebuilding the original environment by downgrading dependencies to the submission date does not help; the reproducible share even falls to 12\,\% \cite{jin2026modernization}.

Repairing such a configuration by hand is slow and demands expertise that many researchers do not have. A line of research has therefore sought to automate it. Classical approaches treat the task as inference over a knowledge base of package metadata \cite{horton2019dockerizeme,ye2022pyego,cheng2022pycre,cheng2024revisiting,wang2021sniffer}. More recent approaches use a \ac{LLM} that reads the error messages of a failed run, proposes a change and retries \cite{bartlett2025lastdependency}, or that acts as an agent which builds a working environment for a whole repository \cite{hu2025repo2run,eliseeva2025envbench}. These systems report encouraging success rates; for example, Repo2Run builds the environment of 86.0\,\% of 420 Python repositories that come with unit tests \cite{hu2025repo2run}.

What these tools have in common is their output. They deliver a corrected configuration, essentially a list of package versions, and the developer must simply trust it. There is no account of \emph{why} a particular version was chosen, whether the choice rests on package-registry facts or on a model's guess, what alternatives were tried and rejected, or whether the result was actually verified. Research on automated program repair suggests that this matters: developers are more willing to accept an automatically generated fix when it comes with explanations and supporting evidence \cite{noller2022trust,ryan2019trust,williams2024bloomberg}, and developers' trust in AI coding tools depends on being able to understand and verify what the tool did \cite{wang2024trust,ye2025copilotlens}. The goal is \emph{calibrated} trust, that is, relying on automation neither too much nor too little \cite{lee2004trust}.% TODO: author to confirm -- the claims in this sentence are summarised from the vision document; the cited papers were checked for existence, authors and venue, but their findings were not re-read for this draft.

\section{Problem Statement}\label{sec:problem}

This thesis addresses three connected problems.

\paragraph{Opaque repairs.}
Existing automated repair of dependency configurations optimises whether the project runs afterwards and treats the reasoning behind each change as disposable. A developer who receives such a repair cannot tell a version that was checked against package metadata from one the model simply suggested, and cannot audit the result.

\paragraph{A mismatch between tools and notebooks.}
The strongest environment-building agents are evaluated on software repositories that ship unit tests, which double as the success signal \cite{hu2025repo2run}, or on single-file code snippets \cite{bartlett2025lastdependency}. Published scientific notebooks have no such tests, and many of their failures are not dependency problems at all, for instance because the data they read were never published. It is unknown which share of real notebook failures a dependency-repair tool can address in principle. Without this figure, a reported success rate cannot be interpreted.

\paragraph{Unverifiable claims of success.}
A repair is only as trustworthy as the evidence behind it. A proposed change should not be accepted because a model says it works, but because an independent check, such as re-running the project, confirms it.

\section{Research Questions}\label{sec:rqs}

The thesis investigates the following research questions. Table~\ref{tab:rqs} shows where each one is answered.

\begin{description}
  \item[RQ1 (Repair).] How effectively can a transparent, agentic \ac{LLM}-based approach repair the dependency configuration of real scientific notebooks so that they execute again?
  \item[RQ2 (Scope).] What share of failing notebooks fails for reasons that dependency repair can address at all, and what are the remaining causes?
  \item[RQ3 (Transparency).] How can the repair process be made transparent, so that each change is recorded with its reason, the source of the decision (package-registry metadata versus model suggestion), its verification status and a resulting confidence level?
  \item[RQ4 (Trust).] Do such explanations help developers to trust correct fixes and to detect incorrect ones? This question is \emph{not} answered in this thesis; it is left as future work (Chapter~\ref{ch:conclusion}).
\end{description}
% TODO: author to confirm -- these four questions follow the advisor-approved outline (RQ1 repair, RQ2 scope, RQ3 transparency, RQ4 trust = future work).
% They differ from the vision document of June 2026, where RQ2 is transparency, RQ3 asks whether transparency costs effectiveness or efficiency
% (compared with a repair-only baseline) and RQ4 is trust. The scope question (new RQ2) emerged from the pilot evaluation. The overhead question is
% not answered by any measurement in the repository and has therefore been dropped here -- confirm this with the advisor.

\begin{table}[ht]
  \centering
  \caption{Research questions and the chapters that address them.}
  \label{tab:rqs}
  \begin{tabular}{@{}l>{\raggedright\arraybackslash}p{0.5\linewidth}>{\raggedright\arraybackslash}p{0.25\linewidth}@{}}
    \hline
    RQ  & Topic & Addressed in \\
    \hline
    RQ1 & Effectiveness of the repair on notebooks & Chapter~\ref{ch:evaluation} \\
    RQ2 & Share of failures that dependency repair can address & Chapters~\ref{ch:background} and~\ref{ch:evaluation} \\
    RQ3 & Transparent, provenance-tagged repair decisions & Chapters~\ref{ch:approach} and~\ref{ch:implementation} \\
    RQ4 & Effect of explanations on developer trust & Future work (Chapter~\ref{ch:conclusion}) \\
    \hline
  \end{tabular}
\end{table}

\section{Contributions}\label{sec:contributions}

The thesis makes the following contributions.

\begin{enumerate}
  \item \textbf{A transparent repair tool for notebooks.} An agentic repair system for Python notebooks and repositories. An \ac{LLM} agent calls a small set of tools to run the project, diagnose the error, look up package metadata, install packages, edit code and verify the result; the model proposes, and only re-execution accepts (Chapter~\ref{ch:approach}).
  \item \textbf{A two-axis provenance model.} Every repair step is labelled along two axes: how the decision was \emph{grounded} (package-registry metadata, deterministic check, or model proposal) and whether it was \emph{verified} by execution. Together, the axes yield a confidence level and a human-readable repair report (Chapter~\ref{ch:approach}).
  \item \textbf{A characterisation of the failures.} A cause-based taxonomy of dependency failures and a measurement of how much of a sample of real notebook failures lies within the reach of dependency repair (Chapters~\ref{ch:background} and~\ref{ch:evaluation}). In a pilot of 26 notebooks, roughly two thirds of the failures that could be examined were not package problems at all.
  \item \textbf{An execution layer and an evaluation harness.} A reproducible runtime for repairs, including a container option, and a harness that records for each notebook why it failed and what the tool did (Chapter~\ref{ch:implementation}). Building it surfaced an engineering finding on Apple-Silicon machines: under the Rosetta emulation used by Docker, a process that holds a ZeroMQ context cannot start any child process, which silently breaks Jupyter-kernel-based tooling.
  \item \textbf{An honest, small-scale evaluation.} Of the 25 notebooks that could be examined, the tool repaired one; this corresponds to 1 of 9--13 genuine package problems, depending on how borderline cases are counted. Six optional repair levers raised this to 2 of 13 on the same 13 in-scope notebooks (four of the 13 were not re-run with the levers). These numbers are too small to support a rate and are reported with their caveats (Chapter~\ref{ch:evaluation}).
\end{enumerate}
% TODO: author to confirm -- numbers in contribution 3 and 5 are taken from PILOT_SUMMARY.md and LEVERS_RESULTS.md in the repository (code/); re-check them against those files after any further runs.

The contribution is deliberately \emph{not} a new agent architecture, nor a record-breaking success rate. Agentic and hybrid designs already exist in recent work, and competing with large agent systems on success rate is not the aim. The contribution is the application to scientific notebooks, the treatment of transparency as a first-class, trust-calibrated output, and the measurement of how much of the problem is in scope for such a tool.% TODO: author to confirm the positioning (framing taken from docs/THESIS_DIRECTION_for_supervisor.md).

The motivation for combining classical and \ac{LLM} techniques comes from the literature itself: Bartlett et al.\ observe that the knowledge-graph baselines still resolve a subset of cases that their \ac{LLM}-based approach fails on, describe the two families as complementary, and name hybrid architectures as future work \cite{bartlett2025lastdependency}. The tool of this thesis uses classical sources of truth, namely package-registry metadata and re-execution, for the facts, and the \ac{LLM} for the reasoning that these sources cannot provide. It does not include a knowledge-graph tool; this is discussed as future work.

\section{Scope and Non-Goals}\label{sec:scope}

The thesis concerns \emph{dependency failures of Python notebooks and scripts}: missing or incompatible packages, interpreter versions, and code that uses an \ac{API} which a newer library version removed. The following are outside its scope.

\begin{itemize}
  \item Failures caused by missing data, hard-coded local paths, external programs or ordinary bugs in the author's own code. The tool detects and reports them, but it does not attempt to fix them (Chapter~\ref{ch:evaluation} quantifies their share).
  \item Software with its own test suite as the success criterion, and the construction of environments for arbitrary repositories at scale.
  \item A study with human participants (RQ4) and a comparison against the tools of the literature on a common benchmark.
  \item Claims of statistical generality. The evaluation is a pilot on a small sample with one language model, one random seed and one machine.
\end{itemize}

\section{Structure of the Thesis}\label{sec:structure}

Chapter~\ref{ch:background} explains what a dependency configuration is, how it decays, why notebooks are especially fragile, and introduces the failure taxonomy. Chapter~\ref{ch:related-work} reviews classical, \ac{LLM}-based and agentic approaches and derives the research gap. Chapter~\ref{ch:approach} presents the concept: the agent loop, its tools, verification by execution and the provenance model. Chapter~\ref{ch:implementation} describes the implementation, including interpreter provisioning, the container runtime and the evaluation harness. Chapter~\ref{ch:evaluation} reports the dataset, the pilot results, the scope finding, the worked examples and the threats to validity. Chapter~\ref{ch:conclusion} summarises the findings and outlines future work. The appendices contain usage notes and the measured data.
